{\scriptsize
        \begin{tabularx}{\textwidth}{|p{0.25\columnwidth}|X|}
            \hline 
            \textbf{Notions et contenus} & \textbf{Capacités exigibles} \\
            \hline 
            \multicolumn{2}{|l|}{\textbf{6.1.3. Ondes électromagnétiques dans le vide }} \\
            \hline
            Équations de propagation d'un champ électromagnétique dans une région sans charge ni courant.  & 
            Établir et citer les équations de propagation d'un champ électromagnétique dans le vide. \\
            \hline 
            Structure d'une onde plane progressive harmonique.   & 
            Établir et exploiter la structure d'une onde électromagnétique plane progressive harmonique.
Utiliser la superposition d'ondes planes progressives harmoniques pour justifier les propriétés d'ondes électromagnétiques planes progressives non harmoniques. \\
        \hline 
         Aspects énergétiques. & 
         Relier la direction du vecteur de Poynting et la direction de propagation de l'onde électromagnétique.
         Interpréter le flux du vecteur de Poynting en termes particulaires.
         Citer quelques ordres de grandeur de flux énergétiques surfaciques moyens et les relier aux ordres de grandeur des champs électriques associés. \\ 
        \hline
        Polarisation des ondes électromagnétiques planes progressives harmoniques : polarisation elliptique, circulaire et rectiligne.  & 
        Relier l'expression du champ électrique à l'état de polarisation de l'onde. \\
        \hline 
        \end{tabularx}
\vspace{0.5cm}
}
